\documentclass[11pt]{article}

\usepackage[letterpaper,margin=1in]{geometry}
\usepackage[T1]{fontenc}
\usepackage[utf8]{inputenc}
\usepackage{lmodern}
\usepackage{microtype}
\usepackage{amsmath,amssymb,mathtools,bm}
\usepackage{enumitem}
\usepackage{xcolor}
\usepackage{booktabs}
\usepackage{array}
\usepackage{hyperref}
\usepackage[most]{tcolorbox}
\usepackage{fancyhdr}
\usepackage{titlesec}
\usepackage{parskip}
\usepackage{needspace}

\definecolor{qftblue}{HTML}{245A8D}
\definecolor{qftgreen}{HTML}{2F6B4F}
\definecolor{qftred}{HTML}{A23B3B}
\definecolor{qftgray}{HTML}{F4F6F8}

\hypersetup{
  colorlinks=true,
  linkcolor=qftblue,
  urlcolor=qftblue,
  citecolor=qftgreen,
  pdftitle={PHYS 662 Lecture 11: Lorentz Symmetry and Its Representations},
  pdfauthor={Nima Lashkari}
}

\pagestyle{fancy}
\fancyhf{}
\lhead{PHYS 662: Quantum Field Theory I}
\rhead{Lecture 11}
\cfoot{\thepage}
\setlength{\headheight}{14pt}

\titleformat{\section}{\Large\bfseries\color{qftblue}}{\thesection.}{0.6em}{}
\titleformat{\subsection}{\large\bfseries\color{qftgreen}}{\thesubsection}{0.6em}{}
\titleformat{\subsubsection}{\normalsize\bfseries\color{qftred}}{\thesubsubsection}{0.6em}{}

\setlist[itemize]{leftmargin=2em,itemsep=0.25em,topsep=0.4em}
\setlist[enumerate]{leftmargin=2em,itemsep=0.25em,topsep=0.4em}

\newcommand{\dd}{\mathrm{d}}
\newcommand{\R}{\mathbb{R}}
\newcommand{\Lag}{\mathcal{L}}
\newcommand{\cM}{\mathcal{M}}
\newcommand{\SO}{\mathrm{SO}}
\newcommand{\OO}{\mathrm{O}}

\newtcolorbox{keyidea}[1][]{
  enhanced, breakable, colback=qftgray, colframe=qftblue,
  boxrule=0.8pt, arc=1.5mm, title={#1}, fonttitle=\bfseries,
  coltitle=white
}
\newtcolorbox{remarkbox}[1][]{
  enhanced, breakable, colback=white, colframe=qftgreen,
  boxrule=0.7pt, arc=1.5mm, title={#1}, fonttitle=\bfseries,
  coltitle=white
}
\newtcolorbox{warningbox}[1][]{
  enhanced, breakable, colback=white, colframe=qftred,
  boxrule=0.7pt, arc=1.5mm, title={#1}, fonttitle=\bfseries,
  coltitle=white
}


\newcommand{\Th}{\Theta}
\newcommand{\Eul}{\mathcal E}
\newcommand{\Lie}{\mathscr L}

\newcommand{\C}{\mathbb C}
\newcommand{\Z}{\mathbb Z}
\newcommand{\SU}{\mathrm{SU}}
\newcommand{\UU}{\mathrm{U}}
\newcommand{\GL}{\mathrm{GL}}
\newcommand{\SL}{\mathrm{SL}}
\newcommand{\Sp}{\mathrm{Sp}}
\newcommand{\tr}{\operatorname{tr}}
\newcommand{\Ad}{\operatorname{Ad}}
\newcommand{\ad}{\operatorname{ad}}
\newcommand{\lieg}{\mathfrak g}
\newcommand{\lieh}{\mathfrak h}
\newcommand{\su}{\mathfrak{su}}
\newcommand{\so}{\mathfrak{so}}
\newcommand{\slie}{\mathfrak{sl}}


\newcommand{\cH}{\mathcal H}
\newcommand{\Spin}{\operatorname{Spin}}
\newcommand{\id}{\mathrm{id}}
\newcommand{\Lor}{\SO^+(1,3)}
\newcommand{\eps}{\varepsilon}
\newcommand{\one}{\mathbf 1}
\begin{document}
\pagenumbering{gobble}
\begin{titlepage}
\centering
\vspace*{1.2in}
{\scshape\Large PHYS 662 --- Quantum Field Theory I\par}
\vspace{0.35in}
{\Huge\bfseries\color{qftblue} Lecture 11\par}
\vspace{0.15in}
{\LARGE\bfseries Lorentz Symmetry\par}
\vspace{0.06in}
{\LARGE\bfseries and Its Representations\par}
\vspace{0.16in}
{\large\color{qftgreen} From Spacetime Isometries to Two Spin Labels\par}
\vspace{0.6in}
{\large Nima Lashkari\par}
{\large Fall 2026\par}
\vfill
\begin{minipage}{0.86\textwidth}
\small
The representation theory of rotations gave us spin. We now include
boosts and ask which kinds of fields can transform covariantly under
Lorentz symmetry. The double cover of the connected Lorentz group and
the complexified Lie algebra lead to two spin labels, Weyl spinors,
and four-vectors. We establish the classification and its limits,
keeping field components, quantum states, and conserved charges distinct.
The explicit spinor calculus begins in Lecture 12.

\medskip
Course page: \url{https://nimalashkari.com/quantum-field-theory-2026-part-i/}
\end{minipage}
\vspace{0.5in}
\end{titlepage}
\pagenumbering{roman}
\begingroup\small
\tableofcontents
\endgroup
\newpage
\pagenumbering{arabic}

\section*{Roadmap}
\addcontentsline{toc}{section}{Roadmap}
Lectures 8--10 developed Lie algebras, spin multiplets, covering groups,
projective phases, and Noether charges. This lecture applies those tools
to spacetime symmetry. The central question is: \emph{what changes when
a spin multiplet must also transform under boosts?}
\begin{enumerate}
\item Describe the Poincar\'e group and identify its connected Lorentz subgroup.
\item Construct the map \(\SL(2,\C)\to\Lor\), including its kernel and topology.
\item Derive the rotation--boost algebra and classify its finite-dimensional
complex representations by \((j_L,j_R)\).
\item Determine rotation content and descent to the Lorentz group,
and distinguish component representations from unitary state representations.
\end{enumerate}
The opening Noether and representation-theory review in the source slides
is already developed in Lectures 8--10. We recall its conclusions here
and concentrate on the Lorentz group and its representations.
Lecture 12 develops Weyl spinors and invariant tensors, Lecture 13 builds
their actions, and Lecture 14 classifies particle states by little groups.

\begin{keyidea}[Conventions]
We use \(c=\hbar=1\), \(\eta_{\mu\nu}=\operatorname{diag}(1,-1,-1,-1)\),
\(\eps_{123}=+1\), and \(x^\mu=(t,\bm x)\). Greek spacetime indices run
from 0 to 3; Latin spatial indices run from 1 to 3. Write
\[
 D(\bm\theta,\bm\rho)=\exp[-i(\theta_iJ_i+\rho_iK_i)]
\]
for a Lorentz representation near the identity. Our active boost along
\(z\) is \(t'=t\cosh\rho-z\sinh\rho\),
\(z'=z\cosh\rho-t\sinh\rho\); reversing \(\rho\) reverses the boost.
The names left and right below refer to the explicitly defined
\((j_L,j_R)\) convention. They are not labels for the sign of a field's momentum.
\end{keyidea}

\section{Spacetime symmetry and the spaces on which it acts}
\subsection{The Poincar\'e group law}
A global isometry of Minkowski space has the form
\begin{equation}
 x'^\mu=\Lambda^\mu{}_{\nu}x^\nu+a^\mu,
 \qquad \Lambda^T\eta\Lambda=\eta.
 \label{eq:isometry}
\end{equation}
Here \(\Lambda\) is a real matrix and \(a\) is a constant displacement.
The isometry condition concerns the metric, not a particular field
configuration. For example, it preserves both spacetime intervals and
\(p_\mu p^\mu=E^2-|\bm p|^2\).

Apply \((\Lambda_2,a_2)\) first and \((\Lambda_1,a_1)\) second:
\[
 x\longmapsto\Lambda_2x+a_2
 \longmapsto\Lambda_1\Lambda_2x+\Lambda_1a_2+a_1.
\]
Therefore
\begin{equation}
 \begin{split}
 (\Lambda_1,a_1)(\Lambda_2,a_2)
   &=(\Lambda_1\Lambda_2,a_1+\Lambda_1a_2),\\
 (\Lambda,a)^{-1}&=(\Lambda^{-1},-\Lambda^{-1}a).
 \end{split}
 \label{eq:group-law}
\end{equation}
The translations form a normal abelian subgroup, but Lorentz
transformations rotate and boost their parameters:
\[
 (\Lambda,0)(I,a)(\Lambda,0)^{-1}=(I,\Lambda a).
\]
This is why the Poincar\'e group is the \emph{semidirect} product
\(\R^{1,3}\rtimes\OO(1,3)\). A direct product would have commuting
translation and Lorentz factors, which this calculation rules out.

\subsection{Field components versus quantum states}
For an active transformation, the old field at \(x\) becomes the new
field at \(x'=\Lambda x+a\). If its components transform by \(D\), then
\begin{equation}
 \phi'_A(x')=D(\Lambda)_A{}^B\phi_B(x),\qquad
 \boxed{\phi'_A(x)=D(\Lambda)_A{}^B
 \phi_B\bigl(\Lambda^{-1}(x-a)\bigr).}
 \label{eq:field-action}
\end{equation}
For spinors we must specify a lift of \(\Lambda\) to the covering group;
we will then write \(D(A)\). The inverse argument is required even for
a scalar, whose component matrix is just 1. The component matrix and the
pullback of the spacetime argument together give the full field transformation.

There are three related but different objects:
\begin{itemize}
\item \(D\) acts on finitely many field components at each spacetime point.
Its boost matrices generally are not unitary.
\item The action on field configurations includes the spacetime argument
and is already infinite-dimensional.
\item A quantum symmetry acts on the Hilbert space of states. The connected
Poincar\'e group is implemented unitarily, possibly projectively before
passing to its cover; its charges are self-adjoint.
\end{itemize}
Thus Wigner's theorem does not require the two-by-two matrices of a
Weyl field to be unitary. It applies to symmetry transformations of
quantum rays. A finite list of field components is not the full space of states.
\subsection{What the Noether review supplies}
For a complex scalar multiplet with internal symmetry
\(\Phi\mapsto e^{-i\epsilon^at_a}\Phi\) in the fundamental representation, Lecture 10 found
\begin{equation}
 \begin{split}
 j_a^\mu&=i\bigl(\Phi^\dagger t_a\partial^\mu\Phi
             -\partial^\mu\Phi^\dagger t_a\Phi\bigr),\qquad
 Q_a=\int\dd^3x\,j_a^0,\\
 \{\Phi,Q_a\}_{\rm PB}&=-it_a\Phi,\qquad
 \{Q_a,Q_b\}_{\rm PB}=f_{ab}{}^cQ_c.
 \end{split}
 \label{eq:review-charges}
\end{equation}
The field equations and vanishing boundary flux give charge conservation.
Quantization supplies \([\widehat Q_a,\widehat Q_b]=if_{ab}{}^c\widehat Q_c\)
under the assumptions stated there. For U(N), the number of generators is
\(N^2\), whereas a fundamental field has only \(N\) components.

This calculation explains how an abstract Lie algebra becomes an algebra
of charges. For spacetime symmetries the same logic applies, but the
field variation also contains the displacement term in
Eq.~\eqref{eq:field-action}. Energy, momentum, angular momentum, and
boosts are the corresponding charges.

Fields, however, are not restricted to the fundamental representation. In many physical contexts, such as non-Abelian gauge theories, fields take values directly in the Lie algebra. To see how these transform, suppose we arrange a set of $N^2$ field components $\phi^a$ into a matrix field using the fundamental representation matrices $(t_a)^i_{\: j}$:
\begin{align}
    \phi^i_{\: j} &= \sum_a \phi^a (t_a)^i_{\: j}\ .
\end{align}
Rather than a simple phase rotation, the finite transformation rule for such a matrix field is given by unitary conjugation:
\begin{align}
    \phi^i_{\: j} &\to U^i_{\: k}\phi^k_{\: l}(U^\dagger)^l_{\: j}\ ,
\end{align}
where $U = e^{-i \epsilon^a t_a}$. The field is acted on from the left by the fundamental representation and from the right by the anti-fundamental representation. 

To find how the individual components $\phi^a$ transform, we expand $U\approx 1 - i \epsilon^a t_a$ and $U^\dagger\approx 1 + i \epsilon^a t_a$. The infinitesimal transformation of the matrix is then
\begin{align}
    \delta \phi^i_{\: j} &= -i \epsilon^a [t_a, \phi]^i_{\: j} = -i \epsilon^a \phi^b [t_a, t_b]^i_{\: j}\ .
\end{align}
Substituting the Lie algebra commutation relation $[t_a, t_b]=i f_{ab}^c t_c$, the variation evaluates to
\begin{align}
    \delta \phi^i_{\: j} &= -i \epsilon^a \phi^b (i f_{ab}^c) (t_c)^i_{\: j} = \epsilon^a f_{ab}^c \phi^b (t_c)^i_{\: j}\ .
\end{align}
By equating this result to the variation of the expansion, $\delta \phi^i_{\: j}=\delta \phi^c (t_c)^i_{\: j}$, we extract the transformation rule for the field components themselves:
\begin{align}
    \delta \phi^c &= \epsilon^a f_{ab}^c \phi^b = -i \epsilon^a (-i f_{ab}^c) \phi^b\ .
\end{align}
Defining the generators of the adjoint representation as $(T_a)^c_{\: b}=-i f_{ab}^c$, we recover $\delta \phi^c=-i \epsilon^a (T_a)^c_{\: b} \phi^b$. This explicitly confirms that while the matrix $\phi^i_{\: j}$ is constructed using fundamental matrices, its components $\phi^a$ transform in the adjoint representation.

\section{The Lorentz group and its double cover}
\subsection{Four connected components}
Taking determinants in Eq.~\eqref{eq:isometry} gives
\((\det\Lambda)^2=1\). It follows from the $00$-component of $\Lambda^T\eta\Lambda=\eta$ that
\[
 (\Lambda^0{}_0)^2-\sum_i(\Lambda^i{}_0)^2=1,
 \qquad |\Lambda^0{}_0|\geq1.
\]
Consequently neither the determinant sign nor the time-orientation sign
can change continuously within the group. In \(3+1\) dimensions the four
components have representatives
\begin{center}
\begin{tabular}{lcc}
\toprule
Representative & \(\det\Lambda\) & \(\operatorname{sign}\Lambda^0{}_0\)\\
\midrule
\(I\) & \(+1\) & \(+1\)\\
\(P=\operatorname{diag}(1,-1,-1,-1)\) & \(-1\) & \(+1\)\\
\(T=\operatorname{diag}(-1,1,1,1)\) & \(-1\) & \(-1\)\\
\(PT=-I\) & \(+1\) & \(-1\)\\
\bottomrule
\end{tabular}
\end{center}
The identity component is the \emph{proper orthochronous}, or restricted,
Lorentz group
\begin{equation}
 \Lor=\{\Lambda:\Lambda^T\eta\Lambda=\eta,\ 
 \det\Lambda=1,\ \Lambda^0{}_0\geq1\}.
 \label{eq:restricted}
\end{equation}
The determinant condition alone does not exclude \(PT\).
We write \(\OO(1,3)\) to match our metric; the notation \(\OO(3,1)\)
in the slides describes an isomorphic group with the opposite signature ordering.

\begin{remarkbox}[Discrete spacetime maps and quantum implementers]
The real matrices \(P\) and \(T\) above describe maps of spacetime.
If a positive-energy quantum theory has time-reversal symmetry, its
Hilbert-space implementer is antiunitary. Conjugating \(i\) then changes
its sign, allowing time evolution to reverse with the energy spectrum
unchanged. Parity is normally unitary. Neither symmetry is guaranteed
by invariance under \(\Lor\).

Charge conjugation changes internal charges; it is not an extra component
of the Lorentz group. The existence of a complex-conjugate representation
does not by itself make charge conjugation a symmetry of the action.
\end{remarkbox}

In any \(d+1\) dimensions with \(d\geq1\), the full \(\OO(1,d)\)
still has four components. For even \(d\), reversing \emph{all} spatial
coordinates has determinant \(+1\); a reflection of just one spatial
coordinate supplies a representative with determinant \(-1\). That is why in odd-dimensional QFT, we often talk about $RT$ (Reflection-Time reversal) operation, instead of $PT$.

\subsection{Encode a four-vector in a Hermitian matrix}
The Pauli matrices form a basis for traceless Hermitian two-by-two matrices.
Adding the identity lets us encode all four real components of \(v\):
\begin{equation}
 V(v)=v^0I+v^i\sigma_i
 =\begin{pmatrix}v^0+v^3&v^1-iv^2\\v^1+iv^2&v^0-v^3\end{pmatrix}.
 \label{eq:vector-matrix}
\end{equation}
This is a real linear bijection, with inverse
\begin{equation}
 v^0=\tfrac12\tr V,\qquad v^i=\tfrac12\tr(\sigma_iV).
 \label{eq:trace-inverse}
\end{equation}
The factor \(1/2\) follows from \(\tr I=2\) and
\(\tr(\sigma_i\sigma_j)=2\delta_{ij}\).
Most importantly,
\begin{equation}
 \det V(v)=(v^0)^2-(v^1)^2-(v^2)^2-(v^3)^2=v_\mu v^\mu.
 \label{eq:det-metric}
\end{equation}
Thus, preservation of the Minkowski quadratic form can be expressed as
preservation of a two-by-two determinant.

For \(A\in\SL(2,\C)\), define
\begin{equation}
 V(v')=AV(v)A^\dagger,\qquad v'=\Lambda(A)v.
 \label{eq:cover-map}
\end{equation}
The result is Hermitian, depends real-linearly on \(v\), and has the
same determinant because \(\det A=1\). Since $\SL(2,\mathbb{C})$ matrices keep the determinant fixed, it follows from (\ref{eq:det-metric}) that the norm of a vector is preserved.
Composition works in the required order:
\[
 A_1(A_2VA_2^\dagger)A_1^\dagger
 =(A_1A_2)V(A_1A_2)^\dagger.
\]
Hence \(A\mapsto\Lambda(A)\) is a group homomorphism.

\subsection{Rotations, boosts, and the kernel}
For \(A_z(\theta)=e^{-i\theta\sigma_3/2}\), the upper off-diagonal
entry of \(AVA^\dagger\) is \(e^{-i\theta}(v^1-iv^2)\). Thus
\[
 v'^1=v^1\cos\theta-v^2\sin\theta,\quad
 v'^2=v^1\sin\theta+v^2\cos\theta,\quad v'^0=v^0,\quad v'^3=v^3.
\]
This recovers the rotation cover of Lecture 9. For a boost take instead
\begin{equation}
 A_z(\rho)=e^{-\rho\sigma_3/2}
   =\begin{pmatrix}e^{-\rho/2}&0\\0&e^{\rho/2}\end{pmatrix}.
 \label{eq:spinor-boost}
\end{equation}
The diagonal entries of \(V\) transform as
\(v'^0+v'^3=e^{-\rho}(v^0+v^3)\) and
\(v'^0-v'^3=e^{\rho}(v^0-v^3)\). Adding and subtracting gives
\begin{equation}
 \begin{pmatrix}v'^0\\v'^3\end{pmatrix}
 =\begin{pmatrix}\cosh\rho&-\sinh\rho\\-\sinh\rho&\cosh\rho\end{pmatrix}
 \begin{pmatrix}v^0\\v^3\end{pmatrix}.
 \label{eq:vector-boost}
\end{equation}
The transverse components stay fixed. Rotating the axis produces any
pure boost. Every proper orthochronous Lorentz transformation is a
boost followed by a rotation: first match its action on the future unit
time vector, then observe that the remaining transformation fixes that
vector and is a spatial rotation. The map is therefore onto \(\Lor\).

If \(A\) is in the kernel, applying Eq.~\eqref{eq:cover-map} to \(V=I\)
gives \(AA^\dagger=I\). Applying it to the Pauli matrices then implies
that \(A\) commutes with every two-by-two matrix. Hence \(A=cI\), and
\(\det A=1\) gives \(c=\pm1\). We have proved
\begin{equation}
 \boxed{\Lor\cong\SL(2,\C)/\{I,-I\}.}
 \label{eq:quotient}
\end{equation}
Merely observing that \(A\) and \(-A\) have the same image is not enough;
the kernel argument also rules out any additional identifications.

\subsection{Why this double cover is universal}
Polar decomposition writes \(A=HU\), where \(U\in\SU(2)\) and
\(H\) is positive Hermitian with determinant one. Uniquely,
\(H=e^{\bm b\cdot\bm\sigma}\) for \(\bm b\in\R^3\).
Consequently, as a manifold,
\begin{equation}
 \SL(2,\C)\simeq\R^3\times S^3.
 \label{eq:polar-topology}
\end{equation}
This is a topological statement, not a direct-product group law.
Both factors are simply connected. The map in Eq.~\eqref{eq:quotient}
is a two-sheeted covering map (its differential is an isomorphism and
its kernel is discrete), so it is the universal cover. Therefore
\begin{equation}
 \pi_1(\Lor)=\Z_2,\qquad
 \Spin^+(1,3)\cong\SL(2,\C).
 \label{eq:fundamental}
\end{equation}
Here the superscript on Spin specifies the cover of the identity
component; it is often suppressed. A \(2\pi\) rotation lifts to a path
from \(I\) to \(-I\), rather than a closed loop.

\begin{remarkbox}[The dimension-dependent exception]
The boost--rotation decomposition gives a deformation retraction of
\(\SO^+(1,d)\) onto \(\SO(d)\): the boost parameters form a contractible
\(\R^d\). Thus its fundamental group is \(\Z_2\) for \(d\geq3\),
but \(\Z\) for \(d=2\). In \(2+1\) dimensions the spin double cover
is not the universal cover. The universal cover has infinitely many
sheets, a fact relevant to arbitrary-spin particle representations and
anyons. In \(1+1\) dimensions the connected boost group is already
isomorphic to \((\R,+)\).
\end{remarkbox}

\section{The Poincar\'e algebra and its charges}
\subsection{Infinitesimal Lorentz transformations}
Put \(\Lambda=I+\omega+O(\omega^2)\) in
\(\Lambda^T\eta\Lambda=\eta\). The linear term is
\begin{equation}
 \omega^T\eta+\eta\omega=0,\qquad
 \omega_{\mu\nu}:=\eta_{\mu\rho}\omega^\rho{}_{\nu}
   =-\omega_{\nu\mu}.
 \label{eq:antisym-omega}
\end{equation}
It is the matrix with its first index \emph{lowered} that is antisymmetric.
There are six independent parameters in four dimensions: three rotations
and three boosts. Adding four translations gives ten Poincar\'e generators.
In \(D\) spacetime dimensions the counts are \(D(D-1)/2\) and
\(D(D+1)/2\), respectively.

Write \(D(\Lambda)=I-\tfrac i2\omega_{\mu\nu}M^{\mu\nu}+O(\omega^2)\),
where \(M^{\mu\nu}=-M^{\nu\mu}\). In the four-vector representation,
\begin{equation}
 (M^{\mu\nu})^\rho{}_{\sigma}
 =i\bigl(\eta^{\mu\rho}\delta^\nu{}_{\sigma}
       -\eta^{\nu\rho}\delta^\mu{}_{\sigma}\bigr).
 \label{eq:vector-generators}
\end{equation}
Substituting this into the infinitesimal transformation recovers
\(D(\Lambda)=I+\omega\). Multiplication of these matrices gives
\begin{equation}
 \begin{split}
 [M^{\mu\nu},M^{\rho\sigma}]
 =i\bigl(&\eta^{\nu\rho}M^{\mu\sigma}
          -\eta^{\mu\rho}M^{\nu\sigma}\\
         &-\eta^{\nu\sigma}M^{\mu\rho}
          +\eta^{\mu\sigma}M^{\nu\rho}\bigr).
 \end{split}
 \label{eq:lorentz-covariant}
\end{equation}
The translation part of the semidirect product adds
\begin{equation}
 [P^\mu,P^\nu]=0,\qquad
 [M^{\mu\nu},P^\rho]
 =i\bigl(\eta^{\nu\rho}P^\mu-\eta^{\mu\rho}P^\nu\bigr).
 \label{eq:translations}
\end{equation}
These are representation-independent algebra relations. In a unitary
Hilbert-space representation their generators are self-adjoint; the
finite matrices in Eq.~\eqref{eq:vector-generators} need not be.

\subsection{Rotations and boosts in three-vector notation}
Define
\begin{equation}
 J_i=\tfrac12\eps_{ijk}M^{jk},\qquad K_i=M^{i0},\qquad H=P^0,
 \qquad p_i:=P^i.
 \label{eq:JK-def}
\end{equation}
The symbol \(p_i\) here denotes Cartesian spatial momentum components;
the covariant four-momentum component is \(P_i=-p_i\).
With these definitions Eqs.~\eqref{eq:lorentz-covariant}--\eqref{eq:translations}
become
\begin{equation}
 \boxed{[J_i,J_j]=i\eps_{ijk}J_k,\quad
 [J_i,K_j]=i\eps_{ijk}K_k,\quad
 [K_i,K_j]=-i\eps_{ijk}J_k.}
 \label{eq:JK-algebra}
\end{equation}
The remaining brackets are
\begin{equation}
 \begin{gathered}
 [J_i,p_j]=i\eps_{ijk}p_k,\qquad [J_i,H]=0,\\
 [K_i,p_j]=i\delta_{ij}H,\qquad [K_i,H]=ip_i,\\
 [p_i,p_j]=0,\qquad [H,p_i]=0.
 \end{gathered}
 \label{eq:poincare-three}
\end{equation}
The minus sign in the boost--boost bracket distinguishes Lorentz
symmetry from four-dimensional Euclidean rotations. For an explicit
check, \(K_1=-i(E_{01}+E_{10})\) and
\(K_2=-i(E_{02}+E_{20})\), where \(E_{ab}\) has a single unit entry.
Then
\[
 [K_1,K_2]=-(E_{12}-E_{21})=-iJ_3,
 \qquad J_3=-i(E_{12}-E_{21}).
\]
Thus two boosts in different directions generate a rotation at second order.

\subsection{Conserved boosts can depend explicitly on time}
For fields with suitable falloff, let \(T^{\mu\nu}\) be a conserved
symmetric stress tensor. A Belinfante improvement includes the intrinsic
spin contribution when starting from a nonsymmetric canonical tensor.
The charges can then be written
\begin{equation}
 P^\mu=\int\dd^3x\,T^{0\mu},\qquad
 M^{\mu\nu}=\int\dd^3x\,(x^\mu T^{0\nu}-x^\nu T^{0\mu}).
 \label{eq:stress-charges}
\end{equation}
Using the unimproved canonical tensor instead requires adding an
explicit spin current to the second expression, as in Lecture 7.
Our boost convention gives
\begin{equation}
 K_i=\int\dd^3x\,x^iT^{00}-tP^i.
 \label{eq:boost-charge}
\end{equation}
Indeed, conservation and symmetry of \(T\), followed by integration by
parts, imply
\[
 \frac{\dd}{\dd t}\int\dd^3x\,x^iT^{00}
 =-\int\dd^3x\,x^i\partial_jT^{j0}
 =\int\dd^3x\,T^{i0}=P^i.
\]
The explicit \(-tP^i\) term cancels this change. Equivalently,
\(\partial_tK_i+i[H,K_i]=-p_i+p_i=0\), even though
\([K_i,H]\ne0\).

\begin{remarkbox}[A one-particle check of the signs]
For a classical free relativistic particle, take
\(H=\sqrt{m^2+\bm p^2}\), \(J_i=\eps_{ijk}x_jp_k\), and
\(K_i=x_iH-tp_i\), with \(\{x_i,p_j\}=\delta_{ij}\).
Since \(\partial H/\partial p_i=p_i/H\), direct differentiation gives
\[
 \{K_i,p_j\}=\delta_{ij}H,\quad
 \{K_i,H\}=p_i,\quad
 \{K_i,K_j\}=x_jp_i-x_ip_j=-\eps_{ijk}J_k.
\]
\end{remarkbox}

\section{Why two spin labels appear}
\subsection{Split the complexified algebra}
The rotation subalgebra is familiar. To incorporate boosts, seek linear
combinations that close by themselves. Define
\begin{equation}
 L_i=\tfrac12(J_i+iK_i),\qquad R_i=\tfrac12(J_i-iK_i).
 \label{eq:LR-def}
\end{equation}
For example, using Eq.~\eqref{eq:JK-algebra},
\begin{align*}
 [L_i,L_j]
 &=\tfrac14\bigl([J_i,J_j]+i[J_i,K_j]+i[K_i,J_j]-[K_i,K_j]\bigr)\\
 &=\tfrac14\bigl(2i\eps_{ijk}J_k-2\eps_{ijk}K_k\bigr)
 =i\eps_{ijk}L_k.
\end{align*}
The other brackets give
\begin{equation}
 \boxed{[L_i,L_j]=i\eps_{ijk}L_k,\quad
 [R_i,R_j]=i\eps_{ijk}R_k,\quad [L_i,R_j]=0.}
 \label{eq:LR-algebra}
\end{equation}
The inverse relations are
\begin{equation}
 J_i=L_i+R_i,\qquad K_i=-i(L_i-R_i).
 \label{eq:JK-inverse}
\end{equation}
Thus the complex representation problem reduces to two commuting copies
of the spin algebra already solved in Lecture 9.

\begin{warningbox}[The real form still matters]
The factors of \(i\) in Eq.~\eqref{eq:LR-def} require complexification.
The precise algebra statement is
\[
 \so(1,3)\otimes_{\R}\C
 \cong\slie(2,\C)\oplus\slie(2,\C).
\]
It does \emph{not} say that the real Lie algebra \(\so(1,3)\) is
\(\su(2)\oplus\su(2)\), or that \(\SL(2,\C)\) is the compact group
\(\SU(2)\times\SU(2)\). The latter group is \(\Spin(4)\), the Euclidean
rotation cover. In a unitary Lorentz representation,
\(J_i^\dagger=J_i\) and \(K_i^\dagger=K_i\) imply \(L_i^\dagger=R_i\);
the two sets are not separately Hermitian spin generators.
\end{warningbox}

\subsection{Construct all finite-dimensional complex irreducibles}
Let \(V_j\) denote the spin-\(j\) representation of the complexified
rotation algebra. Because the two factors commute, each irreducible is
a tensor product
\begin{equation}
 V_{(j_L,j_R)}=V_{j_L}\otimes V_{j_R},\qquad
 j_L,j_R=0,\tfrac12,1,\tfrac32,\ldots,
 \label{eq:irrep}
\end{equation}
with generators
\[
 L_i=J_i^{(j_L)}\otimes I,\qquad
 R_i=I\otimes J_i^{(j_R)}.
\]
One way to see the tensor-product structure is to decompose under the
first factor. The second acts on its multiplicity spaces because it
commutes with the first. Irreducibility of the total action leaves one
first-factor type and an irreducible second-factor multiplicity space.
The dimension is consequently
\begin{equation}
 \boxed{\dim_\C(j_L,j_R)=(2j_L+1)(2j_R+1).}
 \label{eq:dimension}
\end{equation}
These finite-dimensional Lie algebra representations integrate uniquely
to the simply connected real group \(\SL(2,\C)\). The representations
are complex-linear on their vector spaces, but need not depend
holomorphically on \(A\); the second spin label allows dependence on
complex-conjugate matrix entries.

\begin{center}
\begin{tabular}{lll}
\toprule
Lorentz representation & Complex dimension & Example\\
\midrule
\((0,0)\) & 1 & Scalar\\
\((\tfrac12,0)\) & 2 & Left Weyl spinor\\
\((0,\tfrac12)\) & 2 & Right Weyl spinor\\
\((\tfrac12,\tfrac12)\) & 4 & Four-vector\\
\((1,0)\) or \((0,1)\) & 3 each & Chiral antisymmetric tensor\\
\((\tfrac12,0)\oplus(0,\tfrac12)\) & 4 & Dirac spinor (reducible)\\
\bottomrule
\end{tabular}
\end{center}
Dimensions here count complex components before any reality condition.
A real four-vector has four real components, while its complexified
Lorentz representation has complex dimension four. The Dirac spinor is
a \emph{direct sum}; the vector is a \emph{tensor product}.

\subsection{Rotation content and descent to the Lorentz group}
Spatial rotations are generated by \(J_i=L_i+R_i\), so they act
diagonally on the two factors. The usual addition of angular momenta gives
\begin{equation}
 (j_L,j_R)\big|_{\SU(2)_{\rm rot}}
 =\bigoplus_{j=|j_L-j_R|}^{j_L+j_R} V_j,
 \qquad\text{in steps of one.}
 \label{eq:rotation-content}
\end{equation}
For example, \((\tfrac12,\tfrac12)\) restricts to spin zero plus spin
one: its time component is a rotational scalar and its spatial
components form a three-vector. Boosts mix them, so this split is not
a decomposition into Lorentz-invariant subspaces.

The nontrivial kernel element of the cover is a lifted \(2\pi\) rotation.
On the tensor product it acts as
\begin{equation}
 D_{(j_L,j_R)}(-I)=(-1)^{2j_L+2j_R}I.
 \label{eq:kernel-action}
\end{equation}
It follows that
\begin{itemize}
\item \(j_L+j_R\in\Z\): the representation descends to an ordinary
representation of \(\Lor\);
\item \(j_L+j_R\in\Z+\tfrac12\): it is an ordinary representation of
the cover, and only a projective representation of \(\Lor\).
\end{itemize}
For a reducible sum, projective descent requires the kernel to act by
one common scalar on the whole space. The Dirac sum meets this condition
because both summands give \(-I\); a scalar plus a Weyl spinor does not.

\subsection{Topology and central extensions}
The real Lorentz algebra \(\so(1,3)\) is semisimple. Whitehead's lemma,
used in Lecture 10, says that its scalar central Lie algebra extensions
are removable by shifting generators. This is a statement about
two-cocycles, not a consequence of the identity matrix being absent
from a chosen matrix realization. A central extension adds a new
generator outside the original algebra.

The surviving sign in Eq.~\eqref{eq:kernel-action} is global: it comes
from the nontrivial fundamental group. Under the usual continuity and
integrability assumptions for quantum symmetry representations, passing
to \(\SL(2,\C)\) removes this projective ambiguity. For an arbitrary
Lie group, passing to its universal cover need not remove nontrivial
\emph{local} central terms; the Galilean mass example in Lecture 10
already demonstrates this distinction. The Poincar\'e algebra itself
is not semisimple because it contains the translation ideal; its lack
of nontrivial scalar central extensions in \(3+1\) dimensions needs
the separate argument given in Lecture 10.

The nonunitarity of nontrivial finite-dimensional component representations
will be proved alongside the explicit Weyl boosts in Lecture 12.

\section*{Summary}
\addcontentsline{toc}{section}{Summary}
Poincar\'e transformations combine translations and Lorentz transformations
in a semidirect product. The connected Lorentz group has universal cover
\(\SL(2,\C)\), acting on four-vectors through \(V\mapsto AVA^\dagger\).
Its rotation--boost algebra becomes two commuting spin algebras after
complexification. The labels \((j_L,j_R)\) classify finite-dimensional
complex component representations, with dimension
\((2j_L+1)(2j_R+1)\) and rotation spins
\(|j_L-j_R|,\ldots,j_L+j_R\).

The covering kernel determines which representations descend to the
Lorentz group. This global issue is distinct from a local central
extension. Nontrivial finite-dimensional Lorentz representations are
nonunitary, although quantum spacetime symmetries act unitarily on the
infinite-dimensional state space. Lecture 12 turns this classification
into explicit spinor transformations and invariant contractions.

\newpage
\section*{Questions and checks}
\addcontentsline{toc}{section}{Questions and checks}
\begin{enumerate}
\item Derive the Poincar\'e composition and inverse laws by acting on a
point. Verify the field transformation in Eq.~\eqref{eq:field-action}.
\item Show that \(\det\Lambda=1\) does not select the identity component.
Why does \(\OO(1,2)\) still have four components even though reversal of
both spatial coordinates has determinant \(+1\)?
\item Starting from \(V(v)\), derive the trace inverse, determinant,
rotation about \(z\), and boost along \(z\). Prove that the covering
kernel is exactly \(\{I,-I\}\).
\item Explain why polar decomposition makes \(\SL(2,\C)\) simply
connected. Contrast the universal covers in \(3+1\) and \(2+1\) dimensions.
\item Use the four-vector matrices to check
\([J_3,K_1]=iK_2\) and \([K_1,K_2]=-iJ_3\).
Derive \([K_i,H]\) and \([K_i,p_j]\) with \(K_i=M^{i0}\).
\item Verify the boost Poisson brackets in the relativistic-particle
example. How can \(K_i\) be conserved without commuting with \(H\)?
\item Compute \([R_i,R_j]\) and \([L_i,R_j]\). Why do these equations
not identify the real Lorentz group with \(\SU(2)\times\SU(2)\)?
\item Find the dimension, rotation content, and kernel action for
\((1,\tfrac12)\) and \((1,1)\). Which descends to \(\Lor\)?
Does their direct sum define a projective representation on all its rays?
\item Explain why absence of the identity in a matrix algebra does not
prove absence of abstract central extensions. What does the universal
cover remove in this example?
\item Distinguish the spaces on which component matrices, field transformations,
and quantum symmetry operators act. Which of these must be unitary?
\end{enumerate}

\section*{Sources and reading}
\addcontentsline{toc}{section}{Sources and reading}
This lecture develops the group-theoretic material of the supplied
\emph{Lecture 10: Representations of the Lorentz Group} (2025),
especially pp.\ 10--20. Lectures 9 and 10 supply the spin-ladder,
covering-group, central-extension, and Noether-charge prerequisites.
The source's explicit Weyl representations, vector constructions,
epsilon contractions, and statistics discussion continue in Lecture 12;
they have been moved there rather than omitted.
\end{document}
